\chapter{High-Level Design (HLD)}
\section{Architecture Overview}
The system follows a modular event-driven architecture. Each subsystem is isolated by responsibility and communicates through defined interfaces to reduce coupling and improve maintainability.

\section{Major Subsystems}
\begin{itemize}
\item Event Source Layer: generates events from simulation and vision detection.
\item Orchestration Layer: validates events and triggers recording workflow.
\item Recording Layer: builds incident clips from buffered and live frames.
\item Cloud Layer: uploads clips and stores incident metadata.
\item Client Layer: displays incidents and receives notifications.
\end{itemize}

\section{Component Responsibilities}
\begin{longtable}{|p{3.2cm}|p{3.0cm}|p{7.8cm}|}
\hline
\textbf{Component} & \textbf{Primary File/Module} & \textbf{Responsibility} \\
\hline
Event Producer & data\_monitor.py / debug\_update.py & Emit controlled trigger events for test and validation flows. \\
\hline
Event Controller & EventManager.py & Receive and coordinate trigger handling across downstream modules. \\
\hline
Recording Service & recording\_service.py & Manage incident clip creation and recording state transitions. \\
\hline
Cloud/Firebase Adapter & firebase\_manager.py & Push notifications and cloud-facing update operations. \\
\hline
Operator UI & main\_gui.py & Provide local control and status visibility for operators. \\
\hline
Utility Layer & utils.py & Shared helper functions used across modules. \\
\hline
\end{longtable}

\section{Data Flow Summary}
\begin{enumerate}
\item Event source emits trigger.
\item Event controller validates trigger and initiates recording action.
\item Recording service assembles incident clip with contextual buffer.
\item Cloud adapter uploads clip metadata and notifies clients.
\item Client views updated incident list and plays selected clips.
\end{enumerate}

\section{Design Constraints}
\begin{itemize}
\item Timely event processing with low overhead on local hardware.
\item Predictable inter-module interfaces for easier debugging.
\item Compatibility with local backend deployment stack.
\end{itemize}

\section{Design Rationale}
This architecture prioritizes clarity and fault isolation. If one service degrades, logs and boundaries make diagnosis straightforward while preserving the ability to upgrade modules incrementally.